\section{Architecture: Ego-Centric Digital Twin}
\label{sec:architecture}

The architecture has three loops: sensing and aggregate exchange, digital-twin update, and safety-filtered advisory. The real-time loop stays on the vehicle. Offline training may use centralized data and simulation, but deployment does not require roadside sensing, a cloud coordinator, or a globally consistent traffic map.

\noindent\textbf{Sensing and fusion.}
Each vehicle estimates its immediate scene from onboard sensors: ego speed, acceleration, lane, headway, leader and follower gaps, relative speeds, local density, queue estimate, and distance to a merge. Aggregate V2V adds peer-derived summaries over a bounded radius. The combined observation is deliberately not a global state but an ego-local one, with a few semantic hints about what lies outside direct perception. The deployment stack mirrors the simulator's observation contract rather than inventing a separate API, and logs every field with provenance (measured, derived, fallback-neutral, static, or sim-parity). Ego speed, for instance, is measured from GPS, held during short dropouts, and can be cross-checked against on-vehicle odometry; where several sensors observe one quantity the redundancy is reconciled and each field still carries its provenance tag.
Leader gap and closing speed come from forward tracks, density from forward counts, rear-lane fields from neutral fallbacks until a rear camera is added, and nearby-vehicle fields from V2V when enabled. That provenance map is the basis for an observation-quality signal that both the policy and the safety layer can condition on.

\noindent\textbf{The ego-centric digital twin.}
At runtime, each vehicle maintains a lightweight, continuously updated model of its local road state: the ego vehicle, nearby vehicles, lane structure, merge pressure, and short-horizon congestion risk. This is the digital twin. It is a compact tracker, not a city-scale simulator running in the car, and its job is observability. It lets the controller reason about partially observed density and downstream pressure instead of reacting only to the closest leader, and it lets the safety layer test an advisory against a short-horizon prediction before it is shown.

\noindent\textbf{Centralized training and decentralized execution.}
The learning structure is centralized training with decentralized execution (CTDE). During training, a simulator or a logged fleet can expose global state, counterfactual rollouts, and gradients that are unavailable at deployment. During execution, the policy receives only onboard sensing and local-aggregate V2V fields. This split is the technical bridge from the centralized SRC predecessor~\cite{self_regulating_cars} to a vehicle-native system: global information can teach the policy, but it is not part of the runtime control dependency, and no cloud coordinator sits in the loop.

\noindent\textbf{Model-based training for sample efficiency.}
A purely model-free decentralized policy must relearn from many environment interactions each time it meets a new merge geometry, demand pattern, or compliance level. Model-based reinforcement learning is a well-established way to cut that cost: by learning a dynamics model and planning or generating rollouts within it, agents reach comparable performance with substantially fewer real interactions~\cite{sutton1991dyna,janner2019trust,hafner2020dream}. We envision training the self-regulating-car policy in the same way, learning a dynamics model of the road network and using it to adapt the decentralized policy across topologies with few topology-specific interactions. This is a design position: given the sample-efficiency gains consistently reported for model-based RL, a learned road-network model is, in our view, the most promising route to cross-topology adaptation for a policy that must ultimately run on purely local observation.

\noindent\textbf{Advisory control.}
The action space stays conservative: speed bins, headway targets, and cautious lane preferences rather than raw throttle, braking, or aggressive lane-change commands. The public action has four fields: desired speed bin ('slow', 'nominal', 'fast'), desired headway bin ('normal', 'larger', 'largest'), lane preference ('keep', 'prefer left or right if safe'), and merge mode ('normal', 'create gap', 'hold lane'). Prior work on self-regulating cars shows that discrete speed and headway bins of this kind are sufficient for a learned policy to improve flow and reduce hard braking~\cite{self_regulating_cars}. A speed bin is decoded relative to the local segment-target-speed field and floored at a contextual minimum, and the decoded headway is fed back into the next observation. For human-in-the-loop deployment, the output is a bounded advisory in the spirit of prior advisory-autonomy work~\cite{Yan_2023,cho2023temporal,fridman2018human}; for automated vehicles, it can be passed to an existing longitudinal controller. In both cases the self-regulation layer nudges the vehicle toward network-friendly behavior without claiming full authority over the car.

\noindent\textbf{Safety and etiquette filter.}
Every advisory passes through a local filter before it is shown or executed. The filter enforces posted speed limits, acceleration and deceleration bounds, minimum front and rear gaps, target-lane time-to-collision margins, a lane-change dwell time, a cap on lane changes per kilometer, and a limit on the braking an advisory may impose on target-lane followers. It also blocks low-speed behavior in uncongested conditions and synchronized all-lane slowdowns unless downstream congestion justifies them. Blocked or modified actions are surfaced through stable diagnostics: safety-masked action, etiquette block, follower-disruption block, and external override. This layer is not a polish detail; it is the reason a network objective can be exposed safely inside mixed human traffic, and it is what makes the action that is actually executed, rather than the raw policy proposal, the quantity that training and evaluation should measure.

\noindent\textbf{Degraded operation.}
The edge design follows one rule: stale data is worse than missing data. If GPS goes stale, the observation builder holds the last valid speed for a short window rather than reporting a stop; if no V2V peers are heard, cooperation fields revert to explicit neutral values rather than silent zeros; if a display or logger stalls, the pipeline drops outputs rather than queuing an advisory that no longer matches the scene. These modes do not make the system complete, but they turn common in-car failures into explicit states rather than hidden dependencies, and they keep the driver-facing interface advisory-only, surfacing low confidence rather than hiding it inside the policy.
